% !PRINTED 539
\weberpage{539}
%<<CONT
\begin{enumerate}
\item[] and conversely:
\item[2.] \gesperrt{If $\beta_1, \beta_2, \ldots, \beta_n$ is a system of integers
divisible by $\mu$ whose discriminant is equal
to the square of the absolute norm of $\mu$,
multiplied by the fundamental number of the field,
then $\beta_1, \beta_2, \ldots, \beta_n$ is a basis of $\mu$.}
\end{enumerate}

\S.\,148.

\gesperrt{Complete system of residues with respect to a modulus.}

\begin{enumerate}
\item[1.] \textbf{Definition:} \gesperrt{Two integral algebraic numbers
$\xi, \eta$ whose difference $\xi - \eta$ is divisible
by an integral functional $\mu$ different from zero
are called congruent to one another with respect
to the modulus $\mu$.}
\end{enumerate}

The concept of congruence can also be extended to functionals,
which, however, we shall not do for the present. On the other hand,
it is essential to take into account, as moduli of the congruence,
not merely the numbers but also the functionals. Every congruence
remains valid if the modulus is replaced by an associated
functional. In the modulus, arbitrary unit factors may be added
or omitted.

For the congruence we use the \gesperrt{Gauss} symbol
(Bd. I, \S.\,115)
\weq{1}{\xi \equiv \eta \ (\mathrm{mod}\ \mu).}

Such a congruence is equivalent in meaning to the equation
\weq{2}{\xi = \eta + \omega\mu,}
in which $\omega$ can be any integral functional.

From this representation one then recognises at once that, just
as in congruences between rational numbers, if in an expression
composed of integers by \gesperrt{addition}, \gesperrt{subtraction} and \gesperrt{multiplication}
every number is replaced by a number congruent to it with respect
to the modulus $\mu$, the result is a number congruent with respect
to the same modulus; in symbols:

If $\xi_1, \eta_1, \xi_2, \eta_2, \ldots$ are integers which satisfy the congruences
\[
\xi_1 \equiv \eta_1, \ \xi_2 \equiv \eta_2, \ldots \ (\mathrm{mod}\ \mu)
\]
%CONT>>
